%%%%%% My Book Template - preambe
%%%%%% Ali Ashouri
%%%%%%%%%%%%%%%%%%%%%%%%%%%%%%%%%%%%%%%%%%%%%%
%% Packages:
\usepackage[utf8]{inputenc}
\usepackage[a4paper,top=3cm,bottom=2cm,left=0.8in,right=0.8in,marginparwidth=1.75cm]{geometry} %% Sets page size and margins
\usepackage{booktabs,ragged2e,ptext}
\usepackage{amsthm,amssymb,amsmath,amsfonts} %ams
\usepackage{graphicx} %include images/graphics
\usepackage{dsfont} %includes some symbols?
\usepackage{systeme} %system of equations
\usepackage{enumitem} %bullet points
\usepackage{tikz} %draw diagrams
\usepackage{mathtools}
\usepackage{xparse}
\usepackage{bbm} %blackboard fonts
\usepackage{etoolbox} 
\usepackage[dvipsnames]{xcolor}
\usepackage{stmaryrd} %CS font
\usepackage{esvect} %lets me use \vv for a vector arrow
\usepackage[colorlinks=true, allcolors=blue]{hyperref} %makes links blue
\usepackage{thmtools}
\usepackage[framemethod = TikZ]{mdframed}
\usepackage{tikz-cd}
\usepackage{asymptote}
\usepackage{verbatim}
\usepackage[explicit]{titlesec}
\usepackage{epigraph}
\usepackage{float}
%%%%%%%%%%%%%%%%%%%%%%%%%%%%%%%%%%%%%%%%%%%%%%
%% HEADERS AND FOOTERS:
\usepackage{fancyhdr} % Required for customizing headers and footers
\pagestyle{fancy} % Enable the custom headers and footers
\renewcommand{\headrulewidth}{0.5pt} % Top horizontal rule thickness
\fancyhf{} % Clear default headers and footers
\fancyhead[LE, RO]{\sffamily\thepage} % Header for left even pages and right odd pages
\fancyhead[LO]{\rightmark} % Header for left odd pages
\fancyhead[RE]{\leftmark} % Header for right even pages
\fancypagestyle{plain}{ % Style for when a plain pagestyle is specified
	\fancyhead{} % Clear headers
	\renewcommand{\headrulewidth}{0pt} % Remove header rule
}
\usepackage{emptypage} % This package removes headers and footers on empty pages between chapters

%%%%%%%%%%%%%%%%%%%%%%%%%%%%%%%%%%%%%%%%%%%%%%
\usepackage{xepersian} % Persian
%%%%%%%%%%%%%%%%%%%%%%%%%%%%%%%%%%%%%%%%%%%%%%
%% Fonts:
\settextfont{XB Niloofar}
\setlatintextfont{Times New Roman}
%\setdigitfont{Yas}
\settextdigitfont{XB Niloofar}
\setmathdigitfont[Scale=1]{Yas}
\defpersianfont\vazir{Vazirmatn}
\defpersianfont\lalezar{Lalezar-Regular.ttf}
\defpersianfont\yas{Yas}
\defpersianfont\mitra{B Mitra}
\deflatinfont\monotype{Monotype Corsiva}
\deflatinfont\segoe{Segoe UI Black}
%%%%%%%%%%%%%%%%%%%%%%%%%%%%%%%%%%%%%%%%%%%%%%
%% Commands:
\newcommand{\dd}{\,\mathrm{d}}
\newcommand{\se}{\text{} \vspace{-0.45cm} \\}
\newcommand{\see}{\text{} \vspace{-0.75cm}}
%%%%%%%%%%%%%%%%%%%%%%%%%%%%%%%%%%%%%%%%%%%%%%
%% Pages Styles
%\makeatletter
%\def \cleardoublepage{\clearpage \if@twoside \ifodd \c@page \else
%	\hbox{}
%	\thispagestyle{empty}
%	\newpage
%	\if@twocolumn \hbox{} \newpage \fi\fi\fi}
%\makeatother
%%%%%%%%%%%
%% Title Format:

%%%%%%%%%%%%%%%%%%%%%%%%%%%%%%%%%%%%%%%%%%%%%%
%% Environments:
\mdfsetup{skipabove=5pt,skipbelow=0pt} %puts space before and after
\mdfsetup{
	linewidth = 0.3mm,
	innertopmargin = 2mm,
	innerbottommargin = 3.5mm,
	innerleftmargin = 3mm,
	innerrightmargin = 3mm
} % adjusts boundaries of boxes

\newcommand{\thmboxstyle}[4]{
	\mdfdefinestyle{#2}{
		linecolor = #3,
		backgroundcolor = #4,
		nobreak = true
	}
	\declaretheoremstyle[
	headfont = \bfseries\color{#3},
	mdframed = {style = #2},
	headpunct = {\\[0.4pt]},
	postheadspace = {0pt},
	]{#1}
}
\thmboxstyle{defbox}{mdredbox}{red}{orange!5}
\thmboxstyle{thmbox}{mdbluebox}{blue!110!red!90!green!45}{blue!12!red!12!green!1} %purple thmbox
\thmboxstyle{exbox}{mdgreenbox}{blue!75!red!80!green!120}{teal!4!blue!2}
\thmboxstyle{notebox}{mdorangebox}{orange!50!brown}{yellow!5!olive!5}

\declaretheorem[style = thmbox, name = قضیه, numberwithin = chapter]{theorem}
\declaretheorem[style = defbox, name = تعریف, numberwithin = chapter]{definition}
\declaretheorem[style = exbox, name = مثال, numberwithin = chapter]{example}
%%%%%%%%%%%%%%%%%%%%%%%%%%%%%%%%%%%%%%%%%%%%%%

\hypersetup{
	colorlinks=true,
	linkcolor=blue,
	filecolor=magenta,      
	urlcolor=cyan,
	pdftitle={تاریخ جهان در یک نگاه},
	pdfauthor={علی عاشوری},
	pdfsubject={تاریخ جهان},
	pdfkeywords={تاریخ، جهان، تمدن، فرهنگ، علم}
}

%%%%%%%%%%%%%%%%%%%%%%%%%%%%%%%%%%%%%%%%%%%%%%%

% دستورات تاریخ
\newcommand{\bc}[1]{\textbf{#1} ق.م.}  % قبل از میلاد
\newcommand{\ad}[1]{\textbf{#1} م.}     % بعد از میلاد

% دستورات نمایش افراد مهم
\newcommand{\person}[3]{\textbf{#1} (#2--#3)}
\newcommand{\invention}[2]{\textit{#1} (#2)}

% دستورات جعبه اطلاعات
\newcommand{\infobox}[1]{
	\begin{center}
		\fbox{\parbox{0.9\textwidth}{#1}}
	\end{center}
}

% دستور برای خط زمان
\newcommand{\timeline}[3]{
	\noindent
	\begin{minipage}[t]{0.15\textwidth}
		\raggedleft #1
	\end{minipage}%
	\hfill
	\begin{minipage}[t]{0.8\textwidth}
		\textbf{#2}: #3
	\end{minipage}
	\vspace{0.5em}
}

% دستور برای نقل‌قول تاریخی
\newcommand{\historicalquote}[2]{
	\begin{flushright}
		\begin{minipage}{0.8\textwidth}
			\begin{itshape}
				#1
			\end{itshape}
			
			\vspace{0.5em}
			\hfill --- #2
		\end{minipage}
	\end{flushright}
}

%%%%%%%%%%%%%%%%%%%%%%%%%%%%%%%%%%%%%%%
